\subsection{Неравномерные геометрические и амплитудные распределения}\label{sect:distributions}

Антенная решётка - это система из двух или более антенн, которые распределены в пространстве и соединены
электрически таким образом, чтобы сформировать требуемую диаграмму направленности. Электрические
характеристики, направление, поляризация, усиление и взаимное расположение каждого элемента
влияет на итоговые характеристики решетки.

Применение различных неравномерных распределений элементов решётки в пространстве, а также
задание каждому из этих элементов соответствующей амплитуды - это классический подход к оптимизации
антенных решёток. Под оптимизацией понимается уменьшение количества элементов системы, при сохранении
малой ширины основного лепестка и поддержании достаточного ослабления уровня боковых лепестков.
Целью методов, которые применяют такой подход является нахождение подходящего набора положений
и амплитуды элементов антенной решётки при сохранении характеристик, удовлетворяющих
некоторому техническому заданию. Методы этой группы включают в себя использование различных
эмпирических данных, применение интуитивных структур, а также распределений, вдохновлённых
различными природными примерами; применение обратных математических методов, функций оптимизации
и вероятностных подходов, таких как генетические деревья, муравьиный алгоритм и метод роя частиц.

\subsubsection{Интуитивные структуры и эмпирические методы}

Один из применяемых эффективных подходов заключается в использовании интуитивно понятных
геометрических конфигураций в качестве основы, начального  этапа, и проведении измерений
диаграмм направленности при различных случайных или осознанных изменениях в их расположении.
Этот метод позволяет инженерам разрабатывать эффективные антенные решётки без необходимости
полагаться исключительно на сложные вычислительные модели.

\subsubsection{Расширяющееся к краям}

В данном распределении расстояние между элементами тем больше, чем дальше расположен элемент от центра.
Меняя расположение каждого отдельного элемента, можно скорректировать диаграмму направленности
относительно ДН решётки с одинаковым межэлементным расстоянием, уменьшая величину боковых лепестков
при незначительном изменении главного лепестка \cite{harrington1961sidelobe}.

На Рисунке~\ref{fig:sidewiding-distribution-harr} показаны два распределения элементов - равномерное и
расширяющееся к краям; отображена только половина решётки, так как она симметричная.
Неравномерное распределение было рассчитано по методике, описанной в \cite{harrington1961sidelobe}.

\begin{figure}[!ht]
    \centering
    \begin{subfigure}[b]{0.8\textwidth}
        \centering
        \includegraphics[width=\textwidth]{equally-spaced-distribution}
        \caption{}%
        \label{fig:sidewiding-distribution-equal}
    \end{subfigure}
    \begin{subfigure}[b]{0.8\textwidth}
        \centering
        \includegraphics[width=\textwidth]{equally-spaced-distribution}
        \caption{}%
        \label{fig:sidewiding-distribution-unequal}
    \end{subfigure}
    \caption{%
    Расположение элементов 
	(а) равномерное и 
	(б) расширяющееся к краям
    }%
    \label{fig:sidewiding-distribution-harr}
\end{figure}





Программа для моделирования схожа с той, что использовалась в разделе с амплитудными и геометрическими распределениями.

Исходная решётка задавалась следующим образом:


\begin{minted}[
    linenos,
    breaklines,
    frame=single,
    framesep=10pt
]{matlab}
f = 3e9; % 3 GHz text
lam = freq2wavelen(f); % длина волны
xmin = 0;xmax = 29; % количество элементов
dOk = lam/2; % межэлементное расстояние
dxOk = dOk*(xmin:1:xmax)'; % координаты элементов     
\end{minted}

Из неё формируется исследуемая решётка

\begin{minted}[
    linenos,
    breaklines,
    frame=single,
    framesep=10pt
]{matlab}
dxRare = [-2*dOk; dxOk([1 3 5 9:20 26 28 30]); 33*dOk];
dxRare2 = dxOk([9:20]);   
\end{minted}

Далее получаются соответствующие распределения

\begin{minted}[
    linenos,
    breaklines,
    frame=single,
    framesep=10pt
]{matlab}
valuesOk = distribution_former(dxOk,f,aimAngles, aimAmps); % значение сигнала на элементах исходной решётки
valuesRare = distribution_former(dxRare,f,aimAngles, aimAmps);  % значение сигнала на элементах исходной решётки
valuesRare2 = valuesOk([9:20]); % значение сигнала на центральной подрешётке   
\end{minted}

Дополнительно формируются значения со спадающими амплитудными распределениями

\begin{minted}[
    linenos,
    breaklines,
    frame=single,
    framesep=10pt
]{matlab}
valuesRare1 = valuesRare.*([[1,1, 3, 5, 6, 6.5, 7, 7.4, 7.7,8] fliplr([1,1, 3, 5, 6, 6.5, 7, 7.4, 7.7,8])]');
valuesRare21 = valuesRare2.*([[5, 5, 5.5, 6, 6, 6.5] fliplr([5, 5, 5.5, 6, 6, 6.5])]');
\end{minted}

Далее проводится сканирование и перемножение результатов


\begin{minted}[
    linenos,
    breaklines,
    frame=single,
    framesep=10pt
]{matlab}
[rstOk,theta] = aec_simulation(valuesOk, dxOk,f); % ДН эквидистантной решётки для сравнения
rstRare = aec_simulation(valuesRare, dxRare,f); % ДН разрженной решетки для сравнения
rstRare1 = aec_simulation(valuesRare1, dxRare,f);
rstRare2 = aec_simulation(valuesRare2, dxRare2,f); % ДН центарльной решётки для сравнения
rstRare21 = aec_simulation(valuesRare21, dxRare2,f);
rstRare3 = rstRare21.*rstRare1; % результирующая ДН
\end{minted}

Рассмотрим результаты моделирования с одной целью на 0 градусов

\begin{minted}[
    linenos,
    breaklines,
    frame=single,
    framesep=10pt
]{matlab}
aimAngles = [0]; % направления на цели
aimAmps   = [0]; % амплитуды отражённых сигналов на входе приёмника, дБ
\end{minted}